% Part: proof-theory
% Chapter: sequent-calculus
% Section: introduction

\documentclass[../../../include/open-logic-section]{subfiles}

\begin{document}

\olfileid[hi]{pt}{seq}{int}

\olsection{परिचय}

अनुक्रम कलन ऐसी !!{proof} प्रणालियाँ हैं जिन्हें 1930 के दशक में गेर्हार्ड
जेन्ट्सेन ने प्रस्तुत किया था। उनका उद्देश्य वस्तुओं को !!{proving} करने की
व्यावहारिक विधि देना उतना नहीं था। इसके बजाय जेन्ट्सेन ने उनका उपयोग
!!{proof}s की संभावित संरचना और उनके गुणधर्मों के बारे में कुछ परिणाम सिद्ध
करने के लिए किया। प्रमाण सिद्धांत का सरोकार ठीक इसी प्रकार के प्रश्नों से है।

नाम से ही स्पष्ट है कि अनुक्रम कलन !!{formula}s पर नहीं, बल्कि
\emph{अनुक्रमों} पर काम करते हैं। अनुक्रम, !!{formula}s के दो अनुक्रमों या
बहुसमुच्चयों का युग्म होता है। जेन्ट्सेन की मूल प्रणालियों (\Log{LK} और
\Log{LJ}) में अनुक्रमों का उपयोग हुआ था; अब प्रमाण सिद्धांतकार बहुसमुच्चयों
का उपयोग करना पसंद करते हैं। बहुसमुच्चय, समुच्चय जैसा होता है, बस उसमें कोई
!!{element} एक से अधिक बार आ सकता है। फिर भी, समुच्चयों की तरह इसमें भी
!!{element}s का क्रम महत्त्वपूर्ण नहीं होता। अतः व्यंजक $\{!A, !A, !B\}$ और
$\{!A, !B, !A\}$ एक ही बहुसमुच्चय बताते हैं, किंतु वह बहुसमुच्चय $\{!A,
!B\}$ और $\{!A, !A, !B, !B\}$ से भिन्न है, क्योंकि बाद वाले बहुसमुच्चयों में
$!A$ और $!B$ की संख्याएँ पहले वाले से भिन्न हैं।

प्रमाण सिद्धांत में परंपरागत रूप से !!{formula}s के बहुसमुच्चय बड़े यूनानी
अक्षरों से लिखे जाते हैं। इसलिए जब हम $\Gamma$, $\Delta$, $\Pi$, $\Lambda$
या $\Theta$ लिखते हैं, तो उनसे आशय उस तर्क के !!{formula}s के बहुसमुच्चयों
से होता है जिस पर हम काम कर रहे हैं। परंपरागत रूप से प्रमाण सिद्धांतकार
$\tuple{\Gamma, \Delta}$ लिखने के बजाय दोनों घटकों को अनुक्रम-चिह्न से अलग
करते हैं; हम~$\Sequent$ का उपयोग करेंगे।

\begin{defn}[अनुक्रम]
\emph{अनुक्रम} इस रूप का व्यंजक है:
\[
\Gamma \Sequent \Delta
\]
जहाँ $\Gamma$ और $\Delta$, !!{formula}s के सीमित (संभवतः रिक्त) बहुसमुच्चय
हैं। $\Gamma$ को \emph{पूर्वपक्ष} और $\Delta$ को \emph{अनुवर्ती} कहते हैं।
\end{defn}

मान लें कि $!G$, $\Gamma$ के !!{formula}s का संयोजन है (रिक्त संयोजन को
$\ltrue$ मानते हुए), और $!D$, $\Delta$ के !!{formula}s का वियोजन है (रिक्त
वियोजन को $\lfalse$ मानते हुए)। तब अनुक्रम $\Gamma \Sequent \Delta$ (किसी
संरचना~$\Struct{M}$ में) सत्य है यदि और केवल यदि $!G \lif !D$ सत्य है।
चिरसम्मत तर्क में यह कहने के बराबर है कि या तो $\Gamma$ का कोई !!{formula}
असत्य है, अथवा $\Delta$ का कोई !!{formula} सत्य है।

यदि $\Gamma$, !!{formula}s का बहुसमुच्चय है, तो $\Gamma$ में $!A$ जोड़ने से
मिले परिणाम को हम $\Gamma, !A$ (या $!A, \Gamma$) लिखते हैं। यदि $\Delta$ भी
!!{formula}s का बहुसमुच्चय है, तो $\Gamma, \Delta$ दोनों अनुक्रमों का
बहुसमुच्चय-संघ है---यदि $\Gamma$ में $!A$ की \emph{बहुलता}~n है (अर्थात
उसमें $!A$ की $n$ प्रतियाँ हैं) और $\Delta$ में उसकी बहुलता~$m$ है, तो
$\Gamma, \Delta$ में $!A$ की बहुलता~$n+m$ है। यदि अनुक्रम के किसी एक ओर का
बहुसमुच्चय रिक्त है, तो हम प्रायः वहाँ कुछ नहीं लिखते। उदाहरणतः,
$\Sequent !A$ ऐसा अनुक्रम है जिसका पूर्वपक्ष रिक्त है और जिसके अनुवर्ती में
$!A$ की बहुलता~$1$ है।

अनुक्रम कलन में !!^a{proof}, अनुक्रमों का एक वृक्ष होता है, जिसमें सबसे ऊपर
के (अथवा आरंभिक) अनुक्रम विचाराधीन प्रणाली की अभिगृहीतियाँ होते हैं; और यदि
कोई अनुक्रम एक या दो अन्य अनुक्रमों के नीचे है, तो वह प्रणाली के किसी निगमन
नियम द्वारा उनसे सही ढंग से प्राप्त होना चाहिए। हम पहले चिरसम्मत तर्क की दो
भिन्न अनुक्रम प्रणालियों, $\Log{G1c}$ और~$\Log{G3c}$, पर विचार करेंगे। उनकी
अभिगृहीतियाँ और नियम \cref{pt:seq:int:tab:G1c,pt:seq:int:tab:G3c} में दिए
गए हैं। \oliflabeldef{fol:seq::chap}{ (जेन्ट्सेन की मूल प्रणाली \Log{LK} का
वर्णन \olref[fol][seq][]{chap} में है।)}{}

\subfile{rules-G1c}
\subfile{rules-G3c}

इन प्रणालियों में सूक्ष्म अंतर हैं, जिनके कारण इनके प्रमाण-सैद्धांतिक
गुणधर्म भी अलग हैं। \Log{G1c} में $!A \Sequent !A$ अभिगृहीतियाँ हैं, जिनमें
कोई अन्य !!{formula} अनुमत नहीं है। \Log{G3c} में $!C, \Gamma \Sequent
\Delta, !C$ रूप की अभिगृहीतियाँ हैं, जिनके $\Gamma$ और~$\Delta$ में मनमाने
``पार्श्व'' !!{formula}s आ सकते हैं, किंतु दोनों ओर आने वाला !!{formula}~$!C$
\emph{परमाण्विक} होना चाहिए। \Log{G1c} में $\LeftR{\land}$ और
$\RightR{\lor}$ नियमों के दो-दो रूप हैं, जबकि \Log{G3c} में प्रत्येक का केवल
एक रूप है। और \Log{G1c} में ``संकुचन'' (\LeftR{\Contraction},
\RightR{\Contraction}) तथा ``दुर्बलीकरण'' (\LeftR{\Weakening},
\RightR{\Weakening}) नामक अतिरिक्त ``संरचनात्मक नियम'' हैं।

प्रणालियाँ \Log{G1c} और \Log{G3c}, चिरसम्मत तर्क के लिए निर्दोष और पूर्ण
हैं। दूसरे शब्दों में, इन प्रणालियों में !!{provable} प्रत्येक अनुक्रम हर
!!{structure} में सत्य है, और हर !!{structure} में सत्य प्रत्येक अनुक्रम का
!!a{proof} है। अतः किसी अनुक्रम का !!a{proof} है \emph{या नहीं}, इस प्रश्न का
उत्तर तर्क की अन्य विधियों (जैसे निदर्श सिद्धांत की विधियों) से भी दिया जा
सकता है। और चूँकि दोनों प्रणालियाँ ठीक उन्हीं अनुक्रमों को !!{prove} करती हैं
जो हर !!{structure} में सत्य हैं, इसलिए विशेषतः वे समान अनुक्रम सिद्ध करती हैं।

फिर भी, प्रमाण सिद्धांत का सरोकार केवल इससे नहीं है कि कोई वस्तु
\emph{सिद्ध की जा सकती है या नहीं}, बल्कि इससे भी है कि उसे \emph{कैसे}
सिद्ध किया जाता है। प्रमाण सिद्धांतकार ऐसे प्रश्न देखते हैं: ``क्या किसी
नियम से हमेशा बचा जा सकता है?'', ``क्या इस नियम को प्रणाली में जोड़ा जा सकता
है, बिना किसी नए अनुक्रम को सिद्ध करने योग्य बनाए?'', या ``क्या एक प्रणाली
के !!{proof}s को दूसरी प्रणाली के !!{proof}s में बदला जा सकता है?'' यदि ऐसे
प्रश्न का उत्तर ``हाँ'' है, तो इस तथ्य का प्रमाण प्रायः !!{proof}s की जटिलता
पर आगमन से, या समतुल्यतः !!{proof}s की संरचना पर आगमन से दिया जाता है। इससे
केवल तथ्य का प्रमाण ही नहीं मिलता (जैसे, यदि \Log{G1c} कोई अनुक्रम सिद्ध
करता है तो \Log{G3c} भी वही अनुक्रम सिद्ध करता है), बल्कि एक प्रक्रिया भी
मिलती है (जैसे \Log{G1c} के !!{proof}s को \Log{G3c} के !!{proof}s में बदलने
की प्रक्रिया)।

प्रमाण सिद्धांत का एक केंद्रीय परिणाम \emph{कट-विलोपन प्रमेय} है। यह दिखाता
है कि हम कट नियम
\[
\Axiom$\Gamma \fCenter \Delta, !A$
\Axiom$!A, \Pi \fCenter \Lambda$
\RightLabel{\Cut}
\BinaryInf$\Gamma, \Pi \fCenter \Delta, \Lambda$
\DisplayProof
\]
को अनुक्रम प्रणालियों में जोड़ सकते हैं, बिना यह बदले कि कौन-से अनुक्रम
!!{provable} हैं। इसके अतिरिक्त, प्रमाण एक ऐसी विधि देता है जो \Log{G1c}
(या \Log{G3c}) के नियमों के साथ \Cut{} का उपयोग करने वाले किसी भी !!{proof}
को ऐसे प्रमाण में बदल देती है जो केवल \Log{G1c} (या \Log{G3c}) के नियमों का
उपयोग करता है। दूसरे शब्दों में, यह प्रक्रिया \Cut{} नियम का उपयोग करने वाले
!!{proof}s से उस नियम का ``विलोपन'' करती है; इसी से इसका नाम आया है।

\end{document}
